% Version: 28 October 2026. Results and all figure captions rewritten from verified/ (IR_paper commit c96d2e8).
% Introduction, Discussion and Methods unchanged in this revision; see manuscript_planning/ for notes.
% Standard article packages preserve the original single-column 10pt format.
\documentclass[fleqn,10pt]{article}
\usepackage[utf8]{inputenc}
\usepackage[T1]{fontenc}
\usepackage[english]{babel}
\usepackage[letterpaper,left=2cm,right=2cm,top=2.25cm,bottom=2.25cm]{geometry}
\usepackage{times,mathptmx}
\usepackage{amsmath,amssymb,bm,booktabs,graphicx,microtype}
\usepackage[font=small,labelfont={bf,sf},labelsep=period,justification=raggedright,singlelinecheck=false]{caption}
\usepackage{titlesec}
\usepackage[colorlinks=true,allcolors=blue]{hyperref}
\usepackage{fancyhdr,lastpage}
\titleformat{\section}{\large\sffamily\bfseries}{\thesection}{0.5em}{}
\titleformat{\subsection}{\sffamily\bfseries}{\thesubsection}{0.5em}{}
\pagestyle{fancy}
\fancyhf{}
\fancyfoot[R]{\small\thepage/\pageref{LastPage}}
\renewcommand{\headrulewidth}{0pt}
\setlength{\headheight}{13pt}
\setlength{\parindent}{0pt}
\setlength{\parskip}{0.45em}
\setlength{\emergencystretch}{2em}
\allowdisplaybreaks
\newcommand{\figureasset}[1]{%
  \IfFileExists{#1}{\includegraphics[width=\linewidth,height=0.60\textheight,keepaspectratio]{#1}}%
  {\fbox{\parbox{0.94\linewidth}{\centering Figure asset \texttt{\detokenize{#1}} is provided in the companion source archive.}}}%
}
\title{Indirect reciprocity with inertial reputation updating}
\author{Wenqiang Zhu\\[0.5em]
\small School of Artificial Intelligence, Beihang University, Beijing 100191, China\\
\small Key Laboratory of Mathematics, Informatics and Behavioral Semantics,\\
\small Beihang University, Beijing 100191, China\\[0.5em]
\small\textit{Complete author list and correspondence to be added.}}
\date{}
\begin{document}
\raggedbottom
\maketitle
\begin{abstract}
\textit{To be completed following the author's outline.}
\end{abstract}
\section*{Introduction}
\textit{To be completed following the author's outline.}
\section*{Results}
\subsection*{Inertial reputation updating}

We study indirect reciprocity in a well-mixed population playing donation games, in which a donor can pay a cost $c$ to give a recipient a benefit $b>c$. Unconditional cooperators (ALLC) intend to help every recipient, unconditional defectors (ALLD) never intend to help, and discriminators (DISC) intend to help recipients they regard as good. Reputations are private: every individual keeps its own record of every other individual, and an observer judges a donor's realized action by applying a shared social norm to the observer's own view of the recipient. We consider the four second-order norms Scoring (SC), Shunning (SH), Stern Judging (SJ) and Simple Standing (SS), which agree on actions towards recipients regarded as good and differ in how they judge cooperation and defection towards recipients regarded as bad. Intended cooperation fails with probability $e_1$, and an observer records the opposite of the prescribed assessment with probability $e_2$ (Methods).

In the standard model of indirect reciprocity, every new assessment replaces the observer's previous image of the donor: a single negative assessment turns a good image bad, and a single positive assessment turns a bad image good (Fig.~\ref{fig:model}a). We relax this assumption by letting observers require consistent evidence before they reverse an established judgement (Fig.~\ref{fig:model}b). A record of an individual in good standing passes through $k_1$ warning states before it becomes bad, and a record of an individual in bad standing passes through $k_2$ suspicion states before it becomes good. Evidence that supports the established judgement resets the opposing path, so that losing good standing requires $k_1+1$ consecutive negative assessments, and regaining it requires $k_2+1$ consecutive positive assessments, of the same individual by the same observer. We refer to $k_1$ as tolerance depth and to $k_2$ as strictness depth. Because the two paths are independent, inertia is directional: an observer can be slow to condemn, slow to forgive, or both.

When the record is consulted, for choosing an action or for assessing someone else's action, it is read out as a binary image. An individual in warning state $W_m$ is still treated as good with probability $\alpha^m$, and an individual in suspicion state $S_n$ is still treated as bad with probability $\beta^n$. We call $\alpha$ tolerance strength and $\beta$ strictness strength. With full readout, $\alpha=\beta=1$, the established image persists until a path is completed; with $\alpha=\beta=0$, the image follows the latest assessment, which recovers instantaneous updating for any depths. The depths thus determine how much consistent evidence completes a reversal, whereas the readout determines how incomplete evidence affects current treatment. Throughout, we compare three settings: no inertia; full one-step inertia, $(k_1,k_2)=(1,1)$ with $\alpha=\beta=1$; and norm-specific multi-step inertia with full readout, introduced below.

In an environment in which an individual receives positive assessments independently with probability $u$, an inertial record reads that individual as good with a stationary probability $\mathcal R(u)$, which is given in closed form in Methods (Eq.~\eqref{eq:methods-reputation-map}). Without inertia, $\mathcal R(u)=u$. Within the mean-field description, inertia changes reputation dynamics only through the shape of this map, and much of what follows can be read from it. We first analyse reputations at fixed strategy frequencies and then study strategy evolution, assuming that reputations equilibrate faster than strategy frequencies change.

\subsection*{The direction of inertia determines whether private judgements identify defectors}

Private assessment is a central obstacle to indirect reciprocity. When observers hold different views of the same recipient, a discriminator who acts on its own view is condemned by observers who disagree with it, and private and noisy information can undermine cooperation under norms that are stable when information is public~\cite{hilbe2018}. We therefore first asked, at fixed strategy composition, whether inertia allows observers to tell cooperative strategies from defectors, and whether it brings different observers to the same judgement. We simulated finite populations containing equal numbers of ALLC, ALLD and DISC players, in which each DISC observer keeps a separate record of every individual. We quantified recognition by the separation $\Delta g$ between the probability that cooperative strategies (DISC and ALLC) are read as good and the corresponding probability for ALLD. We quantified agreement across observers by the mean disagreement $\bar d$, the probability that two observers read the same individual differently (Methods).

Without inertia, each norm fails in its own way (Fig.~\ref{fig:recognition}a--d, top row). Under Scoring, ALLC players are regarded as good and ALLD players as bad, but observers are split about discriminators, because a discriminator who refuses help to a recipient it regards as bad is condemned. Under Shunning, almost everyone is regarded as bad. Under Stern Judging, every player is regarded as good by about half the observers, and observers disagree at chance level. Under Simple Standing, ALLD players retain substantial trust. Full one-step inertia does not remove these failures (Fig.~\ref{fig:recognition}a--d, middle row). Under Shunning, observers now agree almost perfectly, but only because they regard practically everyone as bad; under Stern Judging, disagreement remains at chance level; and under Scoring, observers remain divided about discriminators. Agreement across observers is therefore not sufficient. What matters is agreement on judgements that distinguish behaviour.

Separation and agreement occur together only when inertia is allocated in a norm-specific direction (Fig.~\ref{fig:recognition}e--h). Using $\Delta g\geq0.8$ and $\bar d\leq0.1$ as operational criteria, the qualifying allocations with the smallest total depth $K=k_1+k_2$ are $(k_1,k_2)=(2,0)$ under Scoring, $(3,0)$ under Shunning, $(4,3)$ under Stern Judging and $(1,2)$ under Simple Standing. Across 30 independent simulations, they yield separations between $0.88$ and $0.92$ with mean disagreement between $0.056$ and $0.081$, whereas full one-step inertia meets both criteria under none of the four norms (Extended Data Fig.~\ref{ed:recognition-selection}a--d). The qualifying regions differ in direction. Under Scoring and Shunning, every qualifying allocation has at least as many warning as suspicion states. Stern Judging requires at least four warning and three suspicion states. Simple Standing requires at least two suspicion states. Direction can matter as much as amount: under Stern Judging, most allocations with more suspicion than warning states (17 of 21 in the scanned grid) invert recognition, so that ALLD players are regarded better than ALLC players. The exact minimal integers depend on the chosen criteria, but the norm-specific direction does not (Extended Data Fig.~\ref{ed:recognition-selection}e--h). We use these four allocations as representative multi-step settings below. They were selected for recognition alone and not for any evolutionary property.

The stationary map $\mathcal R$ shows why direction matters. For the two pure allocations with full readout, it takes a particularly simple form,
\begin{equation}
\mathcal R(u)=1-(1-u)^{k_1+1}\quad\text{for }(k_1,0),
\qquad
\mathcal R(u)=u^{k_2+1}\quad\text{for }(0,k_2).
\label{eq:results-pure-maps}
\end{equation}
A record with only a warning path reads an individual as bad only if the last $k_1+1$ assessments of that individual were all negative, whereas a record with only a suspicion path reads an individual as good only if the last $k_2+1$ assessments were all positive. Tolerance therefore raises every reputation above the assessment probability, and strictness lowers every reputation below it. Whether this helps recognition depends on which strategies the norm assesses unfavourably. Under Scoring, ALLC and ALLD receive positive assessments with fixed probabilities $P^+_{C\mid G}$ and $P^+_{D\mid G}=e_2$, whereas discriminators receive them with probability $e_2+(P^+_{C\mid G}-e_2)\bar r$, which falls when many individuals are regarded as bad. A warning path thus protects discriminators against the condemnation of justified defection. It also raises the reputation of ALLD, but only to $\mathcal R(e_2)=1-(1-e_2)^{k_1+1}\approx(k_1+1)e_2$, which is $0.059$ for $(2,0)$, in agreement with the simulations. Under Simple Standing, by contrast, defection against a recipient regarded as bad is approved, so ALLD players collect positive assessments whenever the recipient is regarded as bad, and the convex map of a suspicion path suppresses their reputation. Under Stern Judging, disagreement about recipients produces erroneous assessments in both directions, and both paths are needed.

Inertia also changes how many reputation states a population can sustain. In a population consisting only of discriminators, the probability $r$ of being regarded as good must satisfy $r=\mathcal R[\Phi_{\mathrm R}(r)]$, where $\Phi_{\mathrm R}$ is the probability of a positive assessment in that population (Methods, Eq.~\eqref{eq:methods-resident-assessment}). With full readout and $k_2\geq1$, $\mathcal R(u)\approx(k_1+1)u^{k_2+1}$ for small $u$. A population whose members regard one another as bad therefore cannot generate the consecutive positive assessments needed to rehabilitate anyone, whereas tolerance keeps a population whose members regard one another as good near the top. Full one-step inertia combines both effects and makes Scoring, Shunning and Stern Judging bistable (Extended Data Fig.~\ref{ed:resident-reputation}a--d). Besides a high-reputation state ($r^*\approx0.995$), these norms sustain a stable low state in which discriminators are regarded as good with probability below $0.001$ (Scoring, Shunning) or $0.5$ (Stern Judging); under Simple Standing the high state is unique. The representative multi-step allocations have a single stable state, with $r^*$ between $0.994$ and $1-4\times10^{-7}$. Scans over the readout parameters and over the depths show that such bistability is widespread (Extended Data Fig.~\ref{ed:reputation-planes}). With full readout it occurs in 60 of the 81 depth pairs with $k_1,k_2\leq8$ under Scoring and under Shunning, in 30 under Stern Judging (exactly those with $k_2\geq k_1\geq1$) and in 3 under Simple Standing. Recovery paths that are long relative to the warning path leave only the low state. Under one-step inertia, the reputation a discriminator population settles into therefore depends on its initial reputation, whereas the representative allocations lead to the same high state from any initial reputation.

Because discriminators help those they regard as good, these reputation states translate directly into cooperation. Without inertia, a population of discriminators is regarded as good, and hence helped, with probability $0.34$ under Scoring, $0.02$ under Shunning, $0.5$ under Stern Judging and $0.83$ under Simple Standing. With the representative multi-step allocations this probability exceeds $0.99$ under all four norms. The next question is whether populations evolve towards such states.

\subsection*{Inertia reorganizes the evolutionary dynamics of cooperation}

To study strategy evolution, we let ALLC, ALLD and DISC evolve according to the replicator dynamics. Payoffs are computed from the self-consistent private reputations reached from initially good records at every strategy composition (Methods). For global comparisons, we follow the dynamics from initial compositions sampled uniformly from the strategy simplex and record the mean long-run cooperation $\bar\rho$, the realized cooperation averaged over these initial conditions. Every trajectory converges either to a state without cooperation or to an attractor with positive cooperation. Hence $\bar\rho\approx\omega_C\bar\rho_C$, where $\omega_C$ is the fraction of initial compositions in the cooperative basin and $\bar\rho_C$ is the cooperation level reached within it.

Inertia changes both which cooperative attractors exist and how they are reached (Fig.~\ref{fig:strategy-dynamics}). Without inertia, a population of discriminators is never stable. Under Shunning and Stern Judging, every trajectory leads to defection. Under Scoring and Simple Standing, cooperation is maintained only in mixtures dominated by unconditional cooperators: a line of neutral equilibria with about 21\% discriminators under Scoring, and a stable mixture with about 78\% ALLC under Simple Standing. Full one-step inertia turns the discriminator population into a strict local attractor under all four norms. With the representative multi-step allocations, discriminators remain a strict attractor under Scoring, Shunning and Stern Judging. Under Simple Standing, a population of discriminators can instead be invaded by ALLC, and the cooperative attractor is a stable mixture of about 80\% discriminators and 20\% unconditional cooperators. In every case, a population of defectors also remains a strict attractor. Inertia does not eliminate defection; it changes the cooperative outcomes and the initial conditions that lead to them.

Averaged over initial compositions, the representative multi-step allocations raise mean cooperation under every norm relative to both no inertia and one-step inertia. Mean cooperation rises from $0.61$ without inertia to $0.87$ under Scoring, from $0$ to $0.69$ under Shunning, from $0$ to $0.70$ under Stern Judging and from $0.54$ to $0.74$ under Simple Standing. Full one-step inertia gives intermediate values under Shunning ($0.08$), Stern Judging ($0.23$) and Simple Standing ($0.70$), but slightly lowers cooperation under Scoring ($0.59$). Two changes contribute. First, the cooperative basin expands: relative to one-step inertia, $\omega_C$ increases from $0.61$ to $0.88$ (Scoring), $0.08$ to $0.71$ (Shunning), $0.23$ to $0.71$ (Stern Judging) and $0.71$ to $0.76$ (Simple Standing), which accounts for most of the increase in $\bar\rho$. Second, the cooperative outcomes themselves change. With inertia, they are sustained by discriminators rather than by majorities of unconditional cooperators, and they cooperate at $0.975$ to $0.980$, close to the maximum of $1-e_1=0.98$ that implementation errors permit. Without inertia, the corresponding levels are $0.961$ (Scoring) and $0.970$ (Simple Standing).

The reputation analysis explains where these basins begin. On the edge of the simplex where only ALLD and DISC are present, the normalized payoff advantage of discriminators over defectors is
\begin{equation}
\frac{\Pi_{\mathrm{DISC}}-\Pi_{\mathrm{ALLD}}}{c(1-e_1)}
=\frac bc\,f_{\mathrm{DISC}}\bigl(r^*_{\mathrm{DISC}}-r^*_{\mathrm{ALLD}}\bigr)-\bar r^*.
\label{eq:results-edge-advantage}
\end{equation}
Discriminators gain because other discriminators help them more than they help defectors, and they pay for the help they give to those they regard as good. Whether they can grow from rarity depends on how their reputation changes as their frequency declines, which inertia reshapes in two distinct ways (Extended Data Fig.~\ref{ed:edge-geometry}). Under Shunning and Stern Judging, one-step inertia sustains a high-reputation branch only when discriminators are very common ($f_{\mathrm{DISC}}\geq0.97$ and $\geq0.77$, respectively). The branch ends at a fold, where the reputation of discriminators collapses and the sign of Eq.~\eqref{eq:results-edge-advantage} reverses abruptly. The boundary of the cooperative basin on this edge is then a reputational tipping point rather than an equilibrium of the strategy dynamics; it produces the abrupt changes in direction visible near the discriminator vertex in the middle row of Fig.~\ref{fig:strategy-dynamics}b,c. The representative allocations remove this fold under Stern Judging, leaving a continuous threshold at $f_{\mathrm{DISC}}\approx0.28$, and displace it under Shunning to $f_{\mathrm{DISC}}\approx0.51$. Under Scoring and Simple Standing, no fold bounds the basin; inertia instead smoothly lowers the frequency at which discriminators begin to outperform defectors, from $0.21$ to $0.08$ under Scoring and from $0.36$ to $0.21$ under Simple Standing. Fold removal and displacement are therefore a key mechanism under Shunning and Stern Judging but not a general one.

Whether a population of discriminators, once established, resists rare unconditional mutants has a compact analytical answer. Let $r_{\mathrm R}^*$ denote the reputation of residents. A rare mutant does not alter the resident assessment environment, and its reputation follows from $\mathcal R$ applied to the positive-assessment probability its actions earn (Methods). When the resident is regarded better than an ALLD mutant and worse than an ALLC mutant, the discriminator population resists both mutants if and only if
\begin{equation}
R_L=\frac{r_{\mathrm R}^{*}}
{r_{\mathrm R}^{*}
-\mathcal R\!\left[
r_{\mathrm R}^{*}P^+_{D\mid G}
+(1-r_{\mathrm R}^{*})P^+_{D\mid B}
\right]}
<
\frac bc
<
\frac{1-r_{\mathrm R}^{*}}
{\mathcal R\!\left[
r_{\mathrm R}^{*}P^+_{C\mid G}
+(1-r_{\mathrm R}^{*})P^+_{C\mid B}
\right]
-r_{\mathrm R}^{*}}=R_U.
\label{eq:canonical-window}
\end{equation}
The lower bound requires defectors to lose more help than they save in costs. The upper bound requires that ALLC's reputational advantage, gained by also helping recipients regarded as bad, not outweigh its additional costs. Without inertia, this window is empty under all four norms, so no benefit-to-cost ratio protects a discriminator population. Full one-step inertia opens a window from $R_L\approx1.00$ to $R_U\approx5.6$--$6.3$. The representative allocations extend $R_U$ to $230$ (Scoring), about $4\times10^3$ (Shunning) and about $2\times10^4$ (Stern Judging), but reduce it to $4.0$ under Simple Standing. Under Simple Standing, the reference ratio $b/c=5$ therefore lies above the window, which is why the cooperative attractor in Fig.~\ref{fig:strategy-dynamics}d is a mixture. Simple Standing nonetheless reaches the highest mean cooperation of its three settings. Local resistance of a discriminator population and mean cooperation are thus distinct properties.

These comparisons are robust to the economic and informational parameters. Across benefit-to-cost ratios from $1.1$ to $10$, the representative allocations yield the highest mean cooperation under every norm (Extended Data Fig.~\ref{ed:benefit-cost}); under Simple Standing, mean cooperation continues to rise beyond $R_U\approx4$, where discriminators lose protection against ALLC. One-step inertia keeps mean cooperation low under Shunning ($\approx0.08$) and below the no-inertia value under Scoring for $b/c\gtrsim3$. Across implementation and assessment errors up to $0.45$, the representative allocations protect discriminators against both mutants over $60\%$ (Scoring), $67\%$ (Shunning) and $31\%$ (Stern Judging) of the error plane, compared with at most $1.6\%$ under one-step inertia and only the error-free corner without inertia (Extended Data Fig.~\ref{ed:error-invasion}). Simple Standing is again an exception, with a smaller protected region under the multi-step than under the one-step allocation. Finally, we asked whether the reputational bistability identified above makes evolutionary outcomes depend on initial reputations (Extended Data Fig.~\ref{ed:resident-reputation}e). Under one-step inertia it does for Shunning and Stern Judging, where cooperation disappears entirely if all records start bad. Under Scoring, the low-reputation state is confined to compositions close to the discriminator vertex, and outcomes are unaffected. The representative allocations are robust across these contrasting initial reputations: outcomes are identical under Scoring, Shunning and Simple Standing, and under Stern Judging mean cooperation changes only from $0.699$ to $0.706$, because a few initial compositions near the ALLC vertex reach cooperation from all-bad records.

\subsection*{The direction of inertia matters more than its amount}

The representative allocations differ between norms, both in total depth and in direction. Are their evolutionary benefits simply a consequence of more inertia? To answer this, we fixed the total depth $K=k_1+k_2$ and compared every allocation between warning and suspicion states (Fig.~\ref{fig:memory-allocation}). The same amount of inertia can produce very different outcomes. With eight intermediate states under Scoring, mean cooperation ranges from almost zero, when all states are suspicion states, to $0.875$, when all are warning states. Pure strictness, $(0,K)$, yields no cooperation under Scoring, Shunning or Stern Judging at any depth, and under Simple Standing cooperation declines with depth, from $0.55$ at $(0,1)$ to $0.06$ at $(0,8)$.

The best allocation at each depth follows a norm-specific pattern that mirrors the recognition results. Under Scoring and Shunning, the most cooperative allocation devotes every state to tolerance. Under Stern Judging, tolerance comes first, and from $K=3$ onwards the best allocations combine both paths. Under Simple Standing, from $K=2$ onwards, one or two warning states combined with a growing recovery path perform best. At their respective total depths, all four representative allocations, which were selected for recognition alone, lie on this cooperation envelope. More inertia is not automatically better. Under Scoring, the envelope peaks at $K=5$ ($\bar\rho=0.888$) and then declines slightly ($0.875$ at $K=8$), and under Stern Judging, $(4,3)$ and $(5,3)$ give the same cooperation. The decline under Scoring is consistent with Eq.~\eqref{eq:results-pure-maps}: each additional warning state raises the reputation of defectors, $\mathcal R(e_2)$, from $0.06$ at $k_1=2$ to $0.17$ at $k_1=8$.

The complete depth plane confirms these patterns (Extended Data Fig.~\ref{ed:depth-evolution}). It also shows that local protection of a discriminator population requires a warning path. Within the scanned range ($k_1,k_2\leq7$) and under all four norms, discriminators resist both unconditional mutants if and only if $k_1\geq2$, or $k_1=1$ with $k_2\leq1$. Protection and cooperation again need not coincide: under Simple Standing with $(1,2)$, mean cooperation is high although discriminators can be invaded by ALLC.

The same asymmetry appears when the direction of inertia is controlled by readout rather than by depth (Extended Data Fig.~\ref{ed:readout-evolution}). Keeping one-step depths fixed, tolerance-only readout, $(\alpha,\beta)=(1,0)$, makes a discriminator population a strict local attractor under all four norms, whereas strictness-only readout, $(0,1)$, does so under none. Across the readout plane, mean cooperation under Scoring, Shunning and Stern Judging is highest at tolerance-only readout and vanishes at strictness-only readout. Under Simple Standing, cooperation remains between $0.54$ and $0.70$ across the entire plane. Whether set by depth or by readout, tolerance of negative evidence is thus the main route to cooperation under Scoring, Shunning and Stern Judging. Strictness helps only under the two norms that approve defection against recipients regarded as bad, Stern Judging and Simple Standing, and under Stern Judging only in combination with tolerance. Tolerance is not required under every norm: without any inertia, Scoring and Simple Standing sustain substantial cooperation through mixtures with unconditional cooperators.

\subsection*{Tolerance invites opportunistic defection unless warning states carry consequences}

Tolerance protects discriminators against negative assessments caused by errors and private disagreement. The same protection, however, might shelter individuals who generate negative assessments deliberately. All mutants considered so far defect either always or never. We therefore asked whether a discriminator population can resist an opportunistic discriminator (O-DISC), which behaves like DISC except that it withholds intended help from a recipient it regards as good with probability $\delta$, independently in every interaction. This family interpolates between DISC ($\delta=0$) and ALLD ($\delta=1$). In a resident discriminator population, a rare O-DISC mutant has the normalized invasion fitness
\begin{equation}
\Lambda(\delta)
=\frac{\Pi_{\mathrm X}-\Pi_{\mathrm R}}{c(1-e_1)}
=\delta r_{\mathrm R}^*
-\frac{b}{c}\bigl[r_{\mathrm R}^*-r_{\mathrm X}^*(\delta)\bigr],
\label{eq:opportunistic-invasion-results}
\end{equation}
where $r_{\mathrm X}^*(\delta)$ is the mutant's stationary reputation. The first term is the saved cost of cooperation and the second the help lost through a worse reputation; the mutant invades if $\Lambda>0$.

Excluding ALLD does not exclude weaker opportunism (Fig.~\ref{fig:opportunism}a). Under Stern Judging with the representative allocation $(4,3)$ and full readout, the resident population is regarded as good with probability $r_{\mathrm R}^*>0.9999$, and ALLD cannot invade. Yet every O-DISC mutant with $0<\delta<0.365$ can. One-step inertia leaves only a narrow window ($0<\delta<0.008$). The vulnerability is not specific to inertia. A rule that assigns a good image if at least one of two sampled assessments is positive~\cite{michelmata2024}, and a bounded quantitative score~\cite{schmid2023}, also exclude ALLD but admit opportunistic mutants up to $\delta\approx0.14$ and $\delta\approx0.38$, respectively. These benchmarks are included to show that tolerant reputation mechanisms in general face this problem, not to rank the mechanisms. Their parameters were not optimized, and the mechanisms are not matched in memory.

The origin of this vulnerability is apparent from the response to weak opportunism. For any reputation mechanism that maps the probability of positive assessments $u$ to a stationary reputation $\mathcal F(u)$ (with $\mathcal F=\mathcal R$ for inertia),
\begin{equation}
\Lambda'(0)
=r_{\mathrm R}^*\left[1-\frac{b}{c}\,\Delta\Phi^* \mathcal F'(u_{\mathrm R}^*)\right],
\label{eq:opportunistic-slope-results}
\end{equation}
where $u_{\mathrm R}^*$ is the residents' positive-assessment probability and $\Delta\Phi^*$ is the difference in positive-assessment probability between intended cooperation and defection in the resident environment (Methods). Weak opportunism pays whenever reputations respond only weakly to an occasional negative assessment near the resident state. Mechanisms that forgive isolated negative assessments have this property by construction. For the aggregation rule, $\mathcal F'(u)=2(1-u)$, which vanishes as $u\to1$. For inertial records with full strictness readout ($\beta=1$), expanding $\mathcal R$ around $u=1$ gives, to leading order in $1-u_{\mathrm R}^*$,
\begin{equation}
\mathcal R'(u_{\mathrm R}^*)\approx(1-\alpha)+(k_1+1)(k_2+1)\bigl(1-u_{\mathrm R}^*\bigr)^{k_1},
\qquad k_1\geq1.
\label{eq:results-marginal-penalty}
\end{equation}
An occasional defection is thus penalized through two channels: the reduced trust in individuals in warning states, $1-\alpha$, and the chance of completing the warning path, which decays geometrically with tolerance depth. With full tolerance readout, $\alpha=1$, only the second channel remains. It is substantial for short warning paths, which is why one-step inertia leaves only a narrow window and why Simple Standing with $(1,2)$ is not invaded. It is negligible for the deeper warning paths of the other representative allocations, which therefore shelter opportunists as effectively as they shelter cooperators against errors.

Reducing $\alpha$ reopens the first channel without changing how the record is updated. For the Stern Judging allocation $(4,3)$, lowering $\alpha$ from $1$ to $0.8$ removes the invasion window entirely, while the resident population remains highly cooperative ($r_{\mathrm R}^*=0.987$; Fig.~\ref{fig:opportunism}a). More generally, all opportunistic mutants are excluded for $\alpha_L<\alpha\leq\alpha_U$, with $\alpha_L\approx0.567$ and $\alpha_U\approx0.804$ (Fig.~\ref{fig:opportunism}b). The two boundaries have different origins. The upper boundary is set by weak opportunism, $\Lambda'(0)=0$, that is, $(b/c)\,\Delta\Phi^*\mathcal R'(u_{\mathrm R}^*)=1$. When the warning path is deep enough for the second term of Eq.~\eqref{eq:results-marginal-penalty} to be negligible and $r_{\mathrm R}^*\approx1$, this gives
\begin{equation}
\alpha_U\approx1-\frac{c}{b\,(1-e_1)(1-2e_2)},
\label{eq:results-alpha-upper}
\end{equation}
which is $0.79$ at the reference parameters. The exact values for the representative allocations under Scoring, Shunning and Stern Judging lie between $0.804$ and $0.807$. The lower boundary is a fold of the resident population: if warning states are trusted too little, $\mathcal R[\Phi_{\mathrm R}(r)]=r$ and $\mathcal R'[\Phi_{\mathrm R}(r)]\,\Phi_{\mathrm R}'(r)=1$ are satisfied simultaneously at $\alpha_L$, and below it the high-reputation state disappears. For the Stern Judging allocation, the resident then has a reputation below $0.5$ and is invadable. Together, these conditions express a trade-off: tolerance must be strong enough to sustain a cooperative population, but weak enough that warning signs carry consequences.

Protective readouts exist under all four norms, at locations that depend on the norm and the allocation (Extended Data Fig.~\ref{ed:opportunism-security}a). At $b/c=5$, the representative allocations are secure for $0\leq\alpha\leq0.807$ under Scoring, $0.591<\alpha\leq0.806$ under Shunning, $0.567<\alpha\leq0.804$ under Stern Judging and over the entire range under Simple Standing. At $\alpha=0.8$, all four are secure, with resident cooperation between $0.978$ and $0.990$. For one-step inertia, the protective intervals under Shunning and Stern Judging are confined to $\alpha>0.85$; at $\alpha=0.8$, these resident populations have already lost their high-reputation state. Because resident reputations do not depend on $b/c$, $\alpha_L$ is fixed, whereas $\alpha_U$ increases with the benefit-to-cost ratio, as in Eq.~\eqref{eq:results-alpha-upper}. A protective interval therefore exists only if $\alpha_U(b/c)>\alpha_L$ (Extended Data Fig.~\ref{ed:opportunism-security}b--e). At $b/c=2$, this condition fails under Shunning and is barely met under Stern Judging. These results concern the specified family of independent opportunistic defection. They do not establish resistance to other forms of exploitation, such as persistent or temporally correlated opportunism, nor do they imply resistance to ALLC: under Simple Standing, the allocation $(1,2)$ is secure against opportunists but invadable by unconditional cooperators.
\section*{Discussion}
\textit{To be completed following the author's outline.}
\section*{Methods}
\label{sec:methods}

\subsection*{Donation game, strategies and private assessment}

We consider a well-mixed population in which a donor can pay a cost $c$ to provide a benefit $b>c$ to a recipient. Defection provides no benefit and incurs no cost. The resident strategy set is $\mathcal S=\{\mathrm{ALLC},\mathrm{ALLD},\mathrm{DISC}\}$. ALLC always intends to cooperate, ALLD always intends to defect, and DISC intends to cooperate if its private image of the recipient is good and to defect otherwise. In the deterministic analyses, $f_I$ denotes the frequency of strategy $I$, with $\sum_I f_I=1$. Donor and recipient roles are sampled without self-interaction.

Reputations are private. A DISC donor chooses its action using its own image of the recipient, whereas an observer assesses the realized action using the observer's image of that recipient. The four second-order norms share positive assessment of cooperation with a good recipient and negative assessment of defection against a good recipient. Their prescriptions for a recipient read as bad are parameterized by $p$ and $q$, the probabilities of positive assessment after cooperation and defection, respectively:
\begin{equation}
\phi(C,G)=1,\qquad \phi(D,G)=0,\qquad
\phi(C,B)=p,\qquad \phi(D,B)=q.
\label{eq:methods-norm}
\end{equation}
Scoring, Shunning, Stern Judging and Simple Standing correspond to $(p,q)=(1,0),(0,0),(0,1)$ and $(1,1)$, respectively.

Intended cooperation fails and becomes defection with probability $e_1$; intended defection is implemented without error. After applying the norm, the observer reverses the prescribed assessment with probability $e_2$. These errors occur independently. Let $P^+_{A\mid R}$ be the probability of a recorded positive assessment conditional on intended action $A\in\{C,D\}$ and the observer's image $R\in\{G,B\}$ of the recipient. Defining $\epsilon=(1-e_1)(1-e_2)+e_1e_2$, we obtain
\begin{equation}
\begin{aligned}
P^+_{C\mid G}&=\epsilon,
& P^+_{D\mid G}&=e_2,\\
P^+_{C\mid B}&=p(\epsilon-e_2)+q(1-\epsilon-e_2)+e_2,
& P^+_{D\mid B}&=q(1-2e_2)+e_2.
\end{aligned}
\label{eq:methods-effective-assessment}
\end{equation}
The action subscript in Eq.~\eqref{eq:methods-effective-assessment} refers to the intended action; implementation error has already been averaged into these probabilities. Unless varied explicitly, $e_1=e_2=0.02$, $c=1$ and $b/c=5$.

\subsection*{Inertial records and their binary readout}

For every ordered observer--target pair $(i,j)$, the observer retains a latent record $L_{ij}$ in the state space
\begin{equation}
\mathcal L=\{G,W_1,\ldots,W_{k_1},B,S_1,\ldots,S_{k_2}\},
\qquad k_1,k_2\in\mathbb Z_{\geq0}.
\label{eq:methods-latent-states}
\end{equation}
The warning states accumulate evidence against established good standing. A negative assessment advances $G\to W_1\to\cdots\to W_{k_1}\to B$, whereas a positive assessment sends every $W_m$ directly to $G$. The suspicion states accumulate evidence in favor of restoring good standing. A positive assessment advances $B\to S_1\to\cdots\to S_{k_2}\to G$, whereas a negative assessment sends every $S_n$ directly to $B$. Positive assessments leave $G$ unchanged and negative assessments leave $B$ unchanged. An absent path is traversed in one update: $G\to B$ after a negative assessment when $k_1=0$, and $B\to G$ after a positive assessment when $k_2=0$. Thus, reversal requires $k_1+1$ consecutive negative assessments or $k_2+1$ consecutive positive assessments. Evidence supporting the established judgment resets the entire opposing path.

Actions and assessments consult a binary image sampled from the latent record. Its good-readout probability is
\begin{equation}
\sigma(\ell)=
\begin{cases}
1,&\ell=G,\\
\alpha^m,&\ell=W_m,\\
0,&\ell=B,\\
1-\beta^n,&\ell=S_n,
\end{cases}
\qquad 0\leq\alpha,\beta\leq1.
\label{eq:methods-readout}
\end{equation}
Every consultation resamples this binary image independently; consultation does not alter $L_{ij}$. Increasing $\alpha$ increases tolerance of negative evidence, whereas increasing $\beta$ increases strictness about rehabilitation. With $\alpha=\beta=1$, warnings are read as good and suspicion states as bad until a path is completed. The one-step model has $k_1=k_2=1$. At these depths, $(\alpha,\beta)=(0,0)$ reproduces instantaneous binary updating, $(1,0)$ gives tolerance alone, $(0,1)$ gives strictness alone, and $(1,1)$ gives full one-step inertia. The depth pair $(0,0)$ is the binary baseline for the depth scans.

\subsection*{Stationary assessment-to-reputation map}

Let $u$ be the probability of a positive recorded assessment in a fixed assessment environment. Write $T^+$ and $T^-$ for the transition matrices corresponding to positive and negative assessments. The latent chain has transition matrix $T(u)=uT^++(1-u)T^-$. Its stationary distribution satisfies $\boldsymbol\pi=\boldsymbol\pi T(u)$ and $\sum_\ell\pi_\ell=1$. Averaging Eq.~\eqref{eq:methods-readout} over this distribution gives the stationary probability of being read as good:
\begin{equation}
\mathcal R(u;\alpha,\beta,k_1,k_2)=
\frac{
u^{k_2+1}\!\left[1+\sum_{m=1}^{k_1}\{\alpha(1-u)\}^{m}\right]
+(1-u)^{k_1+1}\!\sum_{n=1}^{k_2}u^n(1-\beta^n)
}{(1-u)^{k_1}+u^{k_2}-(1-u)^{k_1}u^{k_2}}.
\label{eq:methods-reputation-map}
\end{equation}
Empty sums are zero; the endpoint values are $\mathcal R(0)=0$ and $\mathcal R(1)=1$. For one-step inertia,
\begin{equation}
\mathcal R(u)=
\frac{u^2[1+\alpha(1-u)]+(1-\beta)u(1-u)^2}
{1-u+u^2}.
\label{eq:methods-one-step-map}
\end{equation}
Setting $\alpha=\beta=0$ gives $\mathcal R(u)=u$. The mapping describes stationary readout probabilities under the specified assessment environment; it does not identify the transient binary readout with a Markov process on two states.

\subsection*{Private-reputation self-consistency and branch selection}

For fixed strategy frequencies, $r_I$ denotes the probability that a type-$I$ target is read as good. We use a mean-field closure in which the donor's and observer's private images of a recipient are conditionally independent given the recipient's strategy. Accordingly,
\begin{equation}
\begin{aligned}
\bar r&=\sum_I f_Ir_I,&
h_{GG}&=\sum_I f_Ir_I^2,\\
h_{BB}&=\sum_I f_I(1-r_I)^2,&
h_{GB}&=\sum_I f_Ir_I(1-r_I).
\end{aligned}
\label{eq:methods-private-moments}
\end{equation}
Here $h_{GB}$ is the probability of each ordered disagreement, so the total disagreement probability is $2h_{GB}$. The positive-assessment probabilities generated by the three strategies are
\begin{equation}
\begin{aligned}
\Phi_{\mathrm{ALLC}}&=\bar rP^+_{C\mid G}+(1-\bar r)P^+_{C\mid B},\\
\Phi_{\mathrm{ALLD}}&=\bar rP^+_{D\mid G}+(1-\bar r)P^+_{D\mid B},\\
\Phi_{\mathrm{DISC}}&=h_{GG}P^+_{C\mid G}
+h_{GB}(P^+_{C\mid B}+P^+_{D\mid G})
+h_{BB}P^+_{D\mid B}.
\end{aligned}
\label{eq:methods-positive-assessment}
\end{equation}
Stationary reputations satisfy $r_I^*=\mathcal R(\Phi_I^*)$. The deterministic reputation calculations select a solution from initially good records. For continuous equilibration, we use the auxiliary relaxation
\begin{equation}
\frac{\mathrm dr_I}{\mathrm d\tau}=\mathcal R(\Phi_I)-r_I,
\qquad r_I(0)=1.
\label{eq:methods-reputation-dynamics}
\end{equation}
This equation has the required stationary self-consistency solutions as its fixed points. Its auxiliary time $\tau$ is distinct from the microscopic event clock. The corresponding asynchronous latent-state mean-field system is $\dot{\boldsymbol\pi}_I=\boldsymbol\pi_I[T(\Phi_I)-\mathbf I]$, with $r_I=\boldsymbol\pi_I\boldsymbol\sigma$; its stationary equations give the same mapping within this closure. Equality of stationary equations alone does not establish equality of transient dynamics or attraction basins.

For a homogeneous DISC population,
\begin{equation}
\Phi_R(r)=r^2P^+_{C\mid G}+r(1-r)(P^+_{C\mid B}+P^+_{D\mid G})
+(1-r)^2P^+_{D\mid B}.
\label{eq:methods-resident-assessment}
\end{equation}
We locate the roots of $F(r)=\mathcal R[\Phi_R(r)]-r$ and classify a non-degenerate root as stable or unstable in the scalar relaxation according to the sign of $F'(r)$. The all-good-selected scalar solution is the highest stable root. Extended Data Fig.~\ref{ed:resident-reputation} reports the stationary roots and the stability of this reduced flow. Root finding partitions $[0,1]$ at the extrema of $F$ before applying bracketed scalar root finding, which permits detection of roots near folds. The drift and depth scans use 2,001-point search grids; the one-step readout plane uses a $161\times161$ grid and a 401-point root-search grid, with finer-grid checks. These classifications refer to the scalar closure; they do not establish the stability of the full latent-state system.

\subsection*{Finite-population simulations and recognition statistics}

The quantitative individual-based reputation simulations contain $N=150$ individuals, with 50 individuals of each strategy. The illustrative image matrices use separately saved realizations with $N=90$, comprising 30 individuals of each strategy. Frequencies remain fixed throughout each run. We store the latent records held by the $m=n_{\mathrm{DISC}}$ DISC observers for all $N$ targets, including self-records. Each event draws a donor and a distinct recipient uniformly from the population, then draws an observer independently and uniformly from the DISC individuals. The observer may coincide with the donor or recipient. The donor selects its intended action, implementation error determines the realized action, and the observer applies the norm using a separately sampled private readout of the recipient. Assessment error determines the recorded sign, which updates the observer's latent record of the donor. This simulation retains target-specific variation and correlations between private records rather than imposing the conditional-independence closure.

One sweep consists of $mN$ events, corresponding to one update per retained record on average. All records start in $G$. The recognition depth grid varies $k_1,k_2$ independently over $0,\ldots,6$ at $\alpha=\beta=1$. The initial depth map in Fig.~\ref{fig:recognition} uses ten independent runs per condition; the selection and sensitivity analysis in Extended Data Fig.~\ref{ed:recognition-selection} use 30. The standard protocol runs for 1,000 sweeps and samples every five sweeps from sweep 800 through sweep 1,000, yielding 41 measurements per run. Two slowly relaxing Stern Judging conditions use longer horizons: $(k_1,k_2)=(4,5)$ uses 3,000 sweeps with 2,400 discarded, and $(3,5)$ uses 50,000 sweeps with 40,000 discarded. Their sampling interval and replicate count are unchanged. Condition-specific seeds are generated from base seed 20260918, norm, depths and replicate index. Matrix panels are separately saved illustrative snapshots; the quantitative depth statistics use time-averaged replicate estimates.

Let $s_{ij}=\sigma(L_{ij})$. For target $j$, its mean goodness and disagreement between distinct DISC observers are
\begin{equation}
g_j=\frac{1}{m}\sum_i s_{ij},\qquad
d_j=\frac{1}{\binom m2}\sum_{i<\ell}
\left[s_{ij}(1-s_{\ell j})+(1-s_{ij})s_{\ell j}\right].
\label{eq:methods-finite-estimators}
\end{equation}
These estimators average over stochastic readout exactly. We average over targets of strategy $I$ to obtain $g_I$ and $d_I$, then over sampling times within each replicate. Reputational separation and target-weighted mean disagreement are
\begin{equation}
\Delta g=\frac{g_{\mathrm{DISC}}+g_{\mathrm{ALLC}}}{2}-g_{\mathrm{ALLD}},
\qquad \bar d=\sum_I\frac{n_I}{N}d_I.
\label{eq:methods-recognition-statistics}
\end{equation}
We report arithmetic means across independent replicate means. Where uncertainty is evaluated, pointwise 95\% percentile intervals use 2,000 bootstrap resamples of these replicate means. We assess stationarity by comparing the early and late halves of the sampling window; a drift interval contained within $[-0.01,0.01]$ is a practical small-drift diagnostic, rather than a proof of mixing.

The cutoffs $\Delta g\geq0.8$ and $\bar d\leq0.1$ are operational selection criteria. Among qualifying architectures, we minimize total depth $K=k_1+k_2$, then break ties by greater separation and lower disagreement. This gives $(2,0),(3,0),(4,3)$ and $(1,2)$ for Scoring, Shunning, Stern Judging and Simple Standing. These pairs are selected from reputation formation, independently of evolutionary performance. Extended Data Fig.~\ref{ed:recognition-selection} varies the separation cutoff from 0.70 to 0.90 in steps of 0.05 and the disagreement cutoff from 0.050 to 0.150 in steps of 0.025.

\subsection*{Payoffs, strategy evolution and cooperation}

Reputation formation is assumed to be fast relative to strategy evolution. For each strategy composition, we compute stationary reputations before evaluating expected payoffs:
\begin{equation}
\begin{aligned}
\Pi_{\mathrm{ALLC}}&=(1-e_1)\{b(f_{\mathrm{ALLC}}+f_{\mathrm{DISC}}r_{\mathrm{ALLC}}^*)-c\},\\
\Pi_{\mathrm{ALLD}}&=(1-e_1)b(f_{\mathrm{ALLC}}+f_{\mathrm{DISC}}r_{\mathrm{ALLD}}^*),\\
\Pi_{\mathrm{DISC}}&=(1-e_1)\{b(f_{\mathrm{ALLC}}+f_{\mathrm{DISC}}r_{\mathrm{DISC}}^*)-c\bar r^*\}.
\end{aligned}
\label{eq:methods-payoffs}
\end{equation}
Strategy frequencies obey
\begin{equation}
\frac{\mathrm df_I}{\mathrm dt}=f_I(\Pi_I-\bar\Pi),
\qquad \bar\Pi=\sum_J f_J\Pi_J.
\label{eq:methods-replicator}
\end{equation}
The intended and realized cooperation rates are, respectively,
\begin{equation}
\rho_C^{\mathrm{int}}=f_{\mathrm{ALLC}}+f_{\mathrm{DISC}}\bar r^*,
\qquad \rho_C=(1-e_1)\rho_C^{\mathrm{int}}.
\label{eq:methods-cooperation}
\end{equation}
Thus, an intended cooperation rate approaching one corresponds to realized cooperation approaching $1-e_1$. The simplex equilibrium percentages and the benefit-to-cost and opportunism cooperation panels use intended cooperation. Figure~\ref{fig:memory-allocation} and the readout and depth cooperation maps in Extended Data Figs.~\ref{ed:readout-evolution} and~\ref{ed:depth-evolution} retain the realized-cooperation scale of their scan data.

For global comparisons we average terminal cooperation over initial strategy compositions sampled uniformly with respect to area on the three-strategy simplex. If $\boldsymbol f_{\mathrm{end}}^{(m)}$ is the terminal composition from initial condition $m$, the estimator is
\begin{equation}
\langle\rho_C\rangle_\Delta=
\frac{1}{M}\sum_{m=1}^{M}\rho_C(\boldsymbol f_{\mathrm{end}}^{(m)}),
\qquad
\langle\rho_C^{\mathrm{int}}\rangle_\Delta=
\frac{1}{M}\sum_{m=1}^{M}\rho_C^{\mathrm{int}}(\boldsymbol f_{\mathrm{end}}^{(m)}).
\label{eq:methods-basin-average}
\end{equation}
We use common initial compositions when comparing parameters. The simplex shading denotes the attraction region of the displayed cooperative state. For the total-depth comparison, we enumerate all allocations $k_1+k_2=K$ and define the upper envelope as the largest observed basin-weighted cooperation at each $K$. This is a maximum over the scanned allocations at fixed depth, rather than a global optimum over unrestricted mechanisms.

\subsection*{Numerical strategy calculations}

The supplied simplex panels evaluate the stationary self-consistency map by synchronous iteration of the two private-image moments $\bar r$ and $h_{GG}$, initialized at one. The stopping criterion is the sum of the absolute changes in these moments below $10^{-10}$, with at most 1,000 iterations. The resulting payoff field is evaluated on a rectangular grid clipped to the simplex, with $84\times74$ grid positions before clipping, to draw streamlines. The plotting routine retains terminal iterates when the stationary-map iteration does not converge. Such points occur under Stern Judging and Simple Standing with inertia; their payoff field and associated basin geometry require a converged reputation calculation before being interpreted as stationary dynamics. Vertex stability uses the two rare-strategy growth rates; the displayed filled or open marker at an edge root refers to the eigenvalue tangent to that invariant edge. The transverse invasion eigenvalue is calculated separately. These edge classifications do not establish stability against a strategy absent from that edge. Synchronous iteration is a numerical stationary-map solver, and its convergence is distinct from the stability of Eq.~\eqref{eq:methods-reputation-dynamics}.

The parameter scans integrate the continuous reputation relaxation by fourth-order Runge--Kutta with step 0.1, using residual tolerances of $10^{-9}$ or $10^{-10}$. Slowly relaxing states are continued with step 0.2 for up to 80,000 additional steps. The equilibrated strategy vector field is tabulated on a triangular mesh and interpolated piecewise linearly for midpoint integration of strategy trajectories. Terminal reputation and cooperation values are recomputed without interpolation. The total-depth scan uses a mesh resolution of 23 and 512 common scrambled Sobol initial compositions, midpoint step 0.1, and up to eight blocks of 4,000 steps. The Scoring binary baseline uses its analytical stationary reputations and midpoint step 0.05. Its mean-cooperation stopping rule requires a change below $2\times10^{-5}$ for two successive block comparisons. Three one-step depth points, under Scoring, Shunning and Simple Standing, did not pass this observable-stability gate within the saved horizon; they are finite-time terminal estimates, not confirmed stationary averages.

The benefit-to-cost scan in Extended Data Fig.~\ref{ed:benefit-cost} uses eight independent scrambled Sobol sets of 512 initial compositions, with seeds 2026--2033 and common compositions within each set. Means and standard errors are calculated across the eight set-level means. This uncertainty concerns the distribution of initial strategy compositions; the dynamical model itself is deterministic. The saved calculations use triangular mesh resolutions of 41 for Scoring and 25 for the other norms, midpoint steps 0.05 and 0.1, respectively, and up to 64 integration blocks. Their stopping gate checks changes in both mean and individual intended cooperation and the 99th percentile of the exact terminal strategy residual for successive blocks. Local invasion boundaries are deterministic and do not depend on the initial-condition samples.

\subsection*{Local resistance to ALLC and ALLD}

For a pure-DISC population on the selected resident branch, write $r_R^*$ for the resident reputation. A rare unconditional mutant does not change the resident environment and receives the reputations
\begin{equation}
\begin{aligned}
r_C^*&=\mathcal R\!\left[r_R^*P^+_{C\mid G}+(1-r_R^*)P^+_{C\mid B}\right],\\
r_D^*&=\mathcal R\!\left[r_R^*P^+_{D\mid G}+(1-r_R^*)P^+_{D\mid B}\right].
\end{aligned}
\label{eq:methods-classical-mutants}
\end{equation}
Their invasion growth rates relative to the resident are
\begin{equation}
\begin{aligned}
\lambda_D&=(1-e_1)\{br_D^*-(b-c)r_R^*\},\\
\lambda_C&=(1-e_1)\{br_C^*-c-(b-c)r_R^*\}.
\end{aligned}
\label{eq:methods-classical-invasion}
\end{equation}
When $r_R^*>r_D^*$ and $r_C^*>r_R^*$, strict resistance to both mutants requires
\begin{equation}
R_L=\frac{r_R^*}{r_R^*-r_D^*}<\frac bc<
R_U=\frac{1-r_R^*}{r_C^*-r_R^*}.
\label{eq:methods-classical-thresholds}
\end{equation}
If $r_R^*\leq r_D^*$, no finite benefit-to-cost ratio excludes ALLD. If $r_C^*\leq r_R^*$, ALLC imposes no finite upper threshold. We classify local outcomes directly from the signs of $\lambda_D$ and $\lambda_C$, including cases without a finite threshold. A zero growth rate is a marginal boundary. These criteria refer to the two tested unconditional mutants and are separate from cooperative attraction basins.

\subsection*{Independent opportunistic under-cooperation and benchmark filters}

An opportunistic discriminator, denoted O-DISC, follows DISC except that an intended cooperation is replaced by intended defection independently with probability $\delta$. This deliberate choice occurs before the implementation error. In the pure-DISC resident environment, define
\begin{equation}
\begin{aligned}
\Phi_C^*&=r_R^*P^+_{C\mid G}+(1-r_R^*)P^+_{C\mid B},\\
\Phi_D^*&=r_R^*P^+_{D\mid G}+(1-r_R^*)P^+_{D\mid B},\\
\Delta\Phi^*&=\Phi_C^*-\Phi_D^*,\qquad
u_R^*=\Phi_R(r_R^*).
\end{aligned}
\label{eq:methods-opportunist-assessment}
\end{equation}
Its positive-assessment probability and stationary reputation are
\begin{equation}
u_X(\delta)=u_R^*-\delta r_R^*\Delta\Phi^*,\qquad
r_X^*(\delta)=\mathcal F[u_X(\delta)],
\label{eq:methods-opportunist-reputation}
\end{equation}
where $\mathcal F$ is the reputation filter used by the resident population. The payoff difference gives Eq.~\eqref{eq:opportunistic-invasion-results}. The endpoints are a neutral DISC copy at $\delta=0$ and ALLD at $\delta=1$. Equation~\eqref{eq:opportunistic-slope-results} tests sufficiently small departures from DISC; it does not alone exclude an interior interval of invading mutants.

For inertia, $\mathcal F=\mathcal R$. For the aggregation benchmark, based on Michel-Mata et al.~\cite{michelmata2024}, the filter assigns a good image when at least $Q$ of $M$ independently sampled assessments are positive, giving
\begin{equation}
\mathcal F_A(u)=\sum_{j=Q}^{M}\binom Mj u^j(1-u)^{M-j}.
\label{eq:methods-aggregation-filter}
\end{equation}
The displayed setting is $M=2,Q=1$, so $\mathcal F_A(u)=2u-u^2$. For the quantitative benchmark, based on Schmid et al.~\cite{schmid2023}, the filter uses an integer score $s\in\{-R,\ldots,R\}$, incremented by one after a positive assessment and decremented by one after a negative assessment, with reflection at the endpoints. Scores at least $S$ are read as good. With $t=u/(1-u)$,
\begin{equation}
\mathcal F_Q(u)=\frac{\sum_{j=R+S}^{2R}t^j}{\sum_{j=0}^{2R}t^j},
\qquad 0<u<1.
\label{eq:methods-quantitative-filter}
\end{equation}
The displayed setting is $R=5,S=0$, and endpoint values are obtained by their limits. Each filter is evaluated at its own all-good-selected self-consistent resident reputation under the same norm and errors. The comparisons concern these specified settings, rather than an optimization over aggregation or quantitative parameters.

The displayed invasion curves locate stationary resident roots on a 4,001-point grid, refine sign changes by bracketed root finding, and retain the highest stable root. The exploitation grid combines 4,001 uniformly spaced points on $[0,1]$, 1,001 points on $[0,0.025]$, and 100 logarithmically spaced positive values from $10^{-9}$ to $10^{-3}$; nonzero invasion crossings are refined by root finding. Aggregation and quantitative resident solutions are obtained from $r(0)=1$ by damped stationary-map iteration, $r_{n+1}=0.8r_n+0.2\mathcal F[\Phi_R(r_n)]$, until the step is below $2\times10^{-14}$.

\subsection*{Readout security and resident branch boundaries}

For every readout setting, we recompute the resident reputation and its assessment environment. Define the reputation loss per unit of saved cooperation as
\begin{equation}
\begin{aligned}
g_{\alpha,\beta}(\delta)
&=\frac{r_R^*-\mathcal R\!\left(u_R^*-\delta r_R^*\Delta\Phi^*;\alpha,\beta,k_1,k_2\right)}
{\delta r_R^*},\qquad \delta>0,\\
g_{\alpha,\beta}(0)&=\Delta\Phi^*\mathcal R'(u_R^*).
\end{aligned}
\label{eq:methods-readout-penalty}
\end{equation}
The derivative is with respect to the positive-assessment probability. Exclusion of the entire independently under-cooperating family is equivalent to
\begin{equation}
\frac bc\min_{0\leq\delta\leq1}g_{\alpha,\beta}(\delta)\geq1.
\label{eq:methods-global-security}
\end{equation}
Equality permits a marginal mutant; strict inequality gives strict exclusion. The minimum includes the analytical limit at zero, the endpoint at one and all admissible interior stationary points of the rational loss function. For the security planes, these stationary points are obtained from the real roots of its derivative numerator and checked within the physical interval. The numerical security tolerance is $10^{-10}$.

For the high-reputation branch, the upper readout boundary $\alpha_U$ is defined by Eq.~\eqref{eq:methods-global-security}. A lower boundary $\alpha_L$, when present, is the saddle-node at which that branch disappears. It is located from
\begin{equation}
F(r,\alpha)=0,\qquad \partial_rF(r,\alpha)=0,
\quad F(r,\alpha)=\mathcal R_{\alpha,\beta}[\Phi_R(r)]-r,
\label{eq:methods-resident-fold}
\end{equation}
with the norm, errors, depths and $\beta$ fixed. This fold is distinct from an invasion boundary and is excluded when a strictly stable high-reputation resident is required. We retain only physical resident roots with residual below $5\times10^{-13}$ and reduced-flow slope below $-10^{-9}$ in the refined security analysis; fold solutions are checked separately. The intended resident cooperation shown in these panels is $r_R^*$. Security statements concern O-DISC with independent under-cooperation over $0<\delta\leq1$, together with the separately stated ALLC and ALLD checks; they do not constitute an unrestricted evolutionary-stability claim.

\begin{thebibliography}{9}
\bibitem{hilbe2018}
Hilbe, C., Schmid, L., Tkadlec, J., Chatterjee, K. \& Nowak, M. A.
Indirect reciprocity with private, noisy, and incomplete information.
\textit{Proc. Natl Acad. Sci. USA} \textbf{115}, 12241--12246 (2018).
\url{https://doi.org/10.1073/pnas.1810565115}.
\bibitem{michelmata2024}
Michel-Mata, S., Kawakatsu, M., Sartini, J., Kessinger, T. A., Plotkin, J. B. \& Tarnita, C. E.
The evolution of private reputations in information-abundant landscapes.
\textit{Nature} \textbf{634}, 883--889 (2024).
\url{https://doi.org/10.1038/s41586-024-07977-x}.
\bibitem{schmid2023}
Schmid, L., Ekbatani, F., Hilbe, C. \& Chatterjee, K.
Quantitative assessment can stabilize indirect reciprocity under imperfect information.
\textit{Nature Communications} \textbf{14}, 2086 (2023).
\url{https://doi.org/10.1038/s41467-023-37817-x}.
\end{thebibliography}
\clearpage



\begin{figure}[p]
\centering
\figureasset{figures/Fig1.pdf}
\caption{\textbf{Inertial updating separates how much evidence reverses a judgement from how incomplete evidence is treated.}
\textbf{a}, Instantaneous updating. An observer's image of a donor is replaced by every new assessment: a single negative assessment (pink arrow) turns a good image ($G$) bad ($B$), and a single positive assessment (blue arrow) turns a bad image good.
\textbf{b}, Inertial updating. A record in good standing passes through $k_1$ warning states, $W_1,\dots,W_{k_1}$, on consecutive negative assessments before it becomes $B$ (deterioration path), and a record in bad standing passes through $k_2$ suspicion states, $S_1,\dots,S_{k_2}$, on consecutive positive assessments before it becomes $G$ (recovery path). An assessment that supports the established judgement returns any warning state directly to $G$ and any suspicion state directly to $B$, and leaves $G$ and $B$ unchanged. Losing good standing therefore requires $k_1+1$, and regaining it $k_2+1$, consecutive assessments of the same individual by the same observer. When consulted for an action or an assessment, warning state $W_m$ is read as good with probability $\alpha^m$ (tolerance strength $\alpha$) and suspicion state $S_n$ is read as bad with probability $\beta^n$ (strictness strength $\beta$); labels indicate the readout of the last state of each path. Instantaneous updating is recovered for $k_1=k_2=0$ or $\alpha=\beta=0$.}
\label{fig:model}
\end{figure}
\clearpage

\begin{figure}[p]
\centering
\figureasset{figures/Fig2.pdf}
\caption{\textbf{Norm-specific allocation of inertia lets private observers identify defectors and agree with one another.}
Finite populations with equal numbers of ALLC, ALLD and DISC players, fixed strategies, private records that all start good, and $e_1=e_2=0.02$.
\textbf{a}--\textbf{d}, Private images held by DISC observers (rows) of DISC, ALLC and ALLD players (column blocks) under Scoring (\textbf{a}), Shunning (\textbf{b}), Stern Judging (\textbf{c}) and Simple Standing (\textbf{d}); blue, good; white, bad. Rows of blocks show no inertia (one-step records read out with $\alpha=\beta=0$, equivalent to instantaneous updating), full one-step inertia ($k_1=k_2=1$, $\alpha=\beta=1$) and the representative multi-step allocations $(k_1,k_2)$ indicated above the blocks ($\alpha=\beta=1$). Each block shows one realization after 1{,}000 sweeps in a population of $N=90$.
\textbf{e}--\textbf{h}, Recognition and agreement across tolerance depth $k_1$ and strictness depth $k_2$ ($\alpha=\beta=1$; $N=150$; one simulation per pair, averaged over sweeps 800--1{,}000). Colour shows the separation $\Delta g=\tfrac12(g_{\mathrm{DISC}}+g_{\mathrm{ALLC}})-g_{\mathrm{ALLD}}$, where $g_I$ is the probability that a DISC observer reads a type-$I$ player as good; circle area shows the mean disagreement $\bar d$, the probability that two DISC observers read the same player differently. Black and red outlines mark no inertia, $(0,0)$, and one-step inertia, $(1,1)$. Stars mark the allocation with the smallest total depth satisfying $\Delta g\geq0.8$ and $\bar d\leq0.1$; replicated estimates and the dependence on these criteria are given in Extended Data Fig.~\ref{ed:recognition-selection}.}
\label{fig:recognition}
\end{figure}
\clearpage

\begin{figure}[p]
\centering
\figureasset{figures/Fig3.pdf}
\caption{\textbf{Inertia creates highly cooperative attractors sustained by discriminators, while defection remains an attractor.}
Replicator dynamics of ALLC, ALLD and DISC under Scoring (\textbf{a}), Shunning (\textbf{b}), Stern Judging (\textbf{c}) and Simple Standing (\textbf{d}). Within each panel, the top simplex shows no inertia, the middle simplex full one-step inertia ($k_1=k_2=1$, $\alpha=\beta=1$) and the bottom simplex the representative multi-step allocation ($\alpha=\beta=1$; $(k_1,k_2)=(2,0)$, $(3,0)$, $(4,3)$ and $(1,2)$, respectively). Vertices correspond to monomorphic populations. At every composition, payoffs are computed from self-consistent private reputations reached from initially good records, and streamlines show the resulting dynamics. Filled circles, stable equilibria; open circles, unstable equilibria. The dashed line in \textbf{a} (top) is a line of neutral equilibria. Without inertia, a population of discriminators is never stable; with inertia it is a strict local attractor except under Simple Standing with $(1,2)$, where the cooperative attractor is an ALLC--DISC mixture. The ALLD vertex is stable in every case. Abrupt changes in direction near the DISC vertex under one-step Shunning and Stern Judging arise where the high-reputation branch ends at a fold (Extended Data Fig.~\ref{ed:edge-geometry}). $b/c=5$, $e_1=e_2=0.02$.}
\label{fig:strategy-dynamics}
\end{figure}
\clearpage

\begin{figure}[p]
\centering
\figureasset{figures/Fig4.pdf}
\caption{\textbf{The direction in which inertia is allocated matters more than its total amount.}
Mean long-run cooperation $\bar\rho$ for every allocation of a total depth $K=k_1+k_2\leq8$ between warning ($k_1$) and suspicion ($k_2$) states under Scoring (\textbf{a}), Shunning (\textbf{b}), Stern Judging (\textbf{c}) and Simple Standing (\textbf{d}), with full readout ($\alpha=\beta=1$). $\bar\rho$ is the realized cooperation at the end of the replicator dynamics, averaged over 512 initial strategy compositions sampled uniformly from the simplex; its maximum is $1-e_1=0.98$. Colour shows the fraction $k_1/K$ of states allocated to tolerance (blue, all tolerance; orange, all strictness; $K=0$ is shown in grey); small horizontal offsets separate overlapping points. The black line connects the most cooperative allocation at each $K$, labelled $(k_1,k_2)$. Allocations with the same total depth range from no cooperation to high cooperation. The best allocation is pure tolerance under Scoring and Shunning, mixed under Stern Judging and recovery-weighted under Simple Standing, and the envelope is not monotonic under Scoring. $b/c=5$, $e_1=e_2=0.02$.}
\label{fig:memory-allocation}
\end{figure}
\clearpage

\begin{figure}[p]
\centering
\figureasset{figures/Fig5.pdf}
\caption{\textbf{Tolerant reputations invite opportunistic defection, but a partial readout of warning states restores resistance.}
A rare opportunistic discriminator (O-DISC) in a resident DISC population withholds intended help from a recipient it regards as good with probability $\delta$, independently in every interaction ($\delta=0$, DISC; $\delta=1$, ALLD). Its normalized invasion fitness is $\Lambda(\delta)=\delta r_{\mathrm R}^*-(b/c)[r_{\mathrm R}^*-r_{\mathrm X}^*(\delta)]$ (Eq.~\eqref{eq:opportunistic-invasion-results}); the mutant invades if $\Lambda>0$. Stern Judging, $b/c=5$, $e_1=e_2=0.02$, $\beta=1$.
\textbf{a}, $\Lambda(\delta)$ for full one-step inertia, the multi-step allocation $(k_1,k_2)=(4,3)$ with $\alpha=1$ and $\alpha=0.8$, an aggregation rule (good image if at least one of two sampled assessments is positive) and a quantitative score (integer score bounded at $\pm5$, read as good if non-negative). All five exclude ALLD ($\Lambda(1)<0$), but all except $(4,3)$ with $\alpha=0.8$ admit an interval of invading mutants starting at $\delta=0$; open circles mark its upper end $\delta_c$ ($0.008$, $0.365$, $0.137$ and $0.383$ for one-step, $(4,3)$ with $\alpha=1$, aggregation and quantitative, respectively). Curves are clipped at $\Lambda=-0.15$.
\textbf{b}, Upper end $\delta_c$ of the invasion interval (top; $0$ if none) and resident reputation $r_{\mathrm R}^*$, the intended cooperation of residents (bottom), against the tolerance strength $\alpha$ for $(4,3)$. All O-DISC mutants are excluded for $\alpha_L<\alpha\leq\alpha_U$ (green), with $\alpha_L\approx0.567$ and $\alpha_U\approx0.804$. Above $\alpha_U$, weak opportunism pays ($\Lambda'(0)>0$); at $\alpha_L$, the high-reputation resident state ends at a fold (open circle) and below it the resident occupies a low-reputation state. Residents are the highest strictly stable self-consistent reputation state.}
\label{fig:opportunism}
\end{figure}
\clearpage

\setcounter{figure}{0}
\renewcommand{\figurename}{Extended Data Figure}
\renewcommand{\theHfigure}{ED\arabic{figure}}
% Extended Data figures are ordered by first citation in the Results.
% Asset file names keep the working numbers of verified/ (see comments).

% ED1 (working ED1; verified/figure2_ed1/figures/extended_data_figure1)
\begin{figure}[p]
\centering
\figureasset{figures/ED_Fig1.pdf}
\caption{\textbf{Replicated simulations confirm the representative allocations and show that their direction, not their exact depths, is robust to the recognition criteria.}
\textbf{a}--\textbf{d}, Separation $\Delta g$ against mean disagreement $\bar d$ (defined in Fig.~\ref{fig:recognition}) for all 49 depth pairs $0\leq k_1,k_2\leq6$ under Scoring (\textbf{a}), Shunning (\textbf{b}), Stern Judging (\textbf{c}) and Simple Standing (\textbf{d}); colour gives the total depth $K=k_1+k_2$. Each point is the mean of 30 independent simulations ($N=150$, equal strategy numbers, $\alpha=\beta=1$, records initially good, $e_1=e_2=0.02$). Shading marks $\Delta g\geq0.8$ and $\bar d\leq0.1$. Stars mark the qualifying allocation with minimal $K$ (ties resolved by larger $\Delta g$, then smaller $\bar d$): $(2,0)$, $\Delta g=0.919$, $\bar d=0.064$; $(3,0)$, $0.906$, $0.069$; $(4,3)$, $0.881$, $0.081$; $(1,2)$, $0.917$, $0.056$. Bootstrap 95\% intervals are narrower than $\pm0.001$. Insets show all pairs; crosses mark no inertia.
\textbf{e}--\textbf{h}, Allocation selected under alternative criteria for the maximal disagreement (horizontal) and minimal separation (vertical); cell colour gives its total depth. Gold outlines mark the default criteria; dark outlines enclose criteria that select the same allocation.}
\label{ed:recognition-selection}
\end{figure}
\clearpage

% ED2 (working ED2; verified/figure3_ed2-ed5_ed9/figures/extended_data_figure2)
\begin{figure}[p]
\centering
\figureasset{figures/ED_Fig2.pdf}
\caption{\textbf{One-step inertia makes the reputation of a discriminator population depend on its initial state, whereas the representative allocations do not.}
\textbf{a}--\textbf{d}, Reputation dynamics of a homogeneous DISC population under Scoring (\textbf{a}), Shunning (\textbf{b}), Stern Judging (\textbf{c}) and Simple Standing (\textbf{d}): drift $\dot r=\mathcal R[\Phi_{\mathrm R}(r)]-r$ of the probability $r$ that a discriminator is regarded as good, without inertia (grey), with full one-step inertia (blue) and with the representative multi-step allocation (orange; depths indicated, $\alpha=\beta=1$). Phase lines above each panel show the direction of the drift; filled and open circles mark stable and unstable reputation states. One-step inertia makes Scoring, Shunning and Stern Judging bistable; every representative allocation has a single, high stable state.
\textbf{e}, Mean long-run cooperation $\bar\rho$ when all private records start good (solid bars) or bad (hatched bars) and reputations then follow the strategy dynamics; 512 initial strategy compositions as in Fig.~\ref{fig:memory-allocation}. Outcomes differ under one-step Shunning and Stern Judging, where cooperation disappears from bad initial records, and marginally under multi-step Stern Judging. $b/c=5$ (panel \textbf{e}), $e_1=e_2=0.02$.}
\label{ed:resident-reputation}
\end{figure}
\clearpage

% ED3 (working ED9; verified/figure3_ed2-ed5_ed9/figures/extended_data_figure9)
\begin{figure}[p]
\centering
\figureasset{figures/ED_Fig9.pdf}
\caption{\textbf{Bistability of discriminator reputations is widespread, but the representative allocations avoid it.}
Colour shows the highest stable reputation $r^*$ of a homogeneous DISC population, and hatching marks parameters with two stable reputation states; a hatched cell is therefore not guaranteed to reach the colour shown.
\textbf{a}--\textbf{d}, One-step depths ($k_1=k_2=1$) across tolerance strength $\alpha$ and strictness strength $\beta$ under Scoring (\textbf{a}), Shunning (\textbf{b}), Stern Judging (\textbf{c}) and Simple Standing (\textbf{d}).
\textbf{e}--\textbf{h}, Full readout ($\alpha=\beta=1$) across depths $0\leq k_1,k_2\leq8$. Squares mark no inertia (black) and one-step inertia (blue); stars mark the representative allocations. Bistability occurs in 60 (Scoring), 60 (Shunning), 30 (Stern Judging; exactly where $k_2\geq k_1\geq1$) and 3 (Simple Standing) of the 81 depth pairs, and long recovery paths relative to the warning path leave only a low state. $e_1=e_2=0.02$; reputations do not depend on $b/c$.}
\label{ed:reputation-planes}
\end{figure}
\clearpage

% ED4 (working ED8; verified/ed8_reputation_geometry/figures/extended_data_figure8)
\begin{figure}[p]
\centering
\figureasset{figures/ED_Fig8.pdf}
\caption{\textbf{Inertia moves the point at which discriminators begin to outperform defectors, by removing or displacing reputation folds under Shunning and Stern Judging and by smooth shifts under Scoring and Simple Standing.}
All self-consistent reputation states on the ALLD--DISC edge of the simplex ($f_{\mathrm{ALLC}}=0$) without inertia (grey), with full one-step inertia (blue) and with the representative multi-step allocations (orange; $\alpha=\beta=1$).
\textbf{a}--\textbf{d}, Reputation $r^*_{\mathrm{DISC}}$ of discriminators against their frequency under Scoring (\textbf{a}), Shunning (\textbf{b}), Stern Judging (\textbf{c}) and Simple Standing (\textbf{d}); solid, stable; dashed, unstable; diamonds, folds.
\textbf{e}--\textbf{h}, Normalized payoff advantage of DISC over ALLD, Eq.~\eqref{eq:results-edge-advantage}, on the state reached from initially good records. Circles mark edge equilibria, at which the advantage changes sign continuously; triangles mark sign changes caused by a fold. One-step inertia sustains a high-reputation state under Shunning and Stern Judging only for $f_{\mathrm{DISC}}\geq0.97$ and $\geq0.77$. The representative allocations remove the fold under Stern Judging (edge equilibrium at $f_{\mathrm{DISC}}=0.28$) and displace it to $0.51$ under Shunning, where the reputation curve forms a hysteresis loop. Under Scoring and Simple Standing, the edge equilibrium moves from $0.21$ to $0.08$ and from $0.36$ to $0.21$, respectively; the one-step fold under Scoring lies close to the DISC vertex. $b/c=5$, $e_1=e_2=0.02$.}
\label{ed:edge-geometry}
\end{figure}
\clearpage

% ED5 (working ED4; verified/figure3_ed2-ed5_ed9/figures/extended_data_figure4)
\begin{figure}[p]
\centering
\figureasset{figures/ED_Fig4.pdf}
\caption{\textbf{The representative allocations yield the highest mean cooperation across benefit-to-cost ratios.}
Mean long-run cooperation $\bar\rho$ against $b/c$ under Scoring (\textbf{a}), Shunning (\textbf{b}), Stern Judging (\textbf{c}) and Simple Standing (\textbf{d}) without inertia (grey), with full one-step inertia (blue) and with the representative multi-step allocations (orange; $\alpha=\beta=1$). Lines interpolate means over eight independent samples of 512 initial compositions at 45 ratios from $1.001$ to $10$; shading, $\pm1$ s.e.m. across samples. Solid segments indicate that a DISC population resists both rare ALLD and rare ALLC ($R_L<b/c<R_U$, Eq.~\eqref{eq:canonical-window}); dashed segments indicate that it does not. Labels give thresholds within the plotted range and arrows indicate $R_U>10$; without inertia, the protected range is empty. Under Simple Standing, the representative allocation loses protection against ALLC above $R_U\approx4.0$, yet its mean cooperation continues to increase. $e_1=e_2=0.02$.}
\label{ed:benefit-cost}
\end{figure}
\clearpage

% ED6 (working ED5; verified/figure3_ed2-ed5_ed9/figures/extended_data_figure5)
\begin{figure}[p]
\centering
\figureasset{figures/ED_Fig5.pdf}
\caption{\textbf{The representative allocations protect discriminator populations against unconditional mutants over a wide range of errors under Scoring, Shunning and Stern Judging.}
Invasion of a resident DISC population by rare ALLD and ALLC across implementation error $e_1$ and assessment error $e_2$ ($241\times241$ grid) under Scoring, Shunning, Stern Judging and Simple Standing (columns), without inertia (\textbf{a}--\textbf{d}), with full one-step inertia (\textbf{e}--\textbf{h}) and with the representative multi-step allocations (\textbf{i}--\textbf{l}; $\alpha=\beta=1$). Pink, ALLD is not excluded (including neutral cases); purple, ALLD is excluded but ALLC invades; green, the resident resists both. Dashed and solid white curves mark zero invasion fitness of ALLD and ALLC, respectively. Residents are reached from initially good records; on the line $e_2=0$, under Scoring and Shunning, this is the all-bad state, against which ALLD is neutral. Green covers 60\%, 67\% and 31\% of the plane under Scoring, Shunning and Stern Judging with multi-step inertia, at most 1.6\% with one-step inertia, and only the error-free corner without inertia. Simple Standing is an exception. $b/c=5$.}
\label{ed:error-invasion}
\end{figure}
\clearpage

% ED7 (working ED6; verified/figure4_ed6/figures/extended_data_figure6)
\begin{figure}[p]
\centering
\figureasset{figures/ED_Fig6.pdf}
\caption{\textbf{Across the full depth plane, high cooperation and local protection of discriminators both require a warning path, but they do not coincide.}
Mean long-run cooperation $\bar\rho$ (colour; computed as in Fig.~\ref{fig:memory-allocation}) for all depth pairs $0\leq k_1,k_2\leq7$ under Scoring (\textbf{a}), Shunning (\textbf{b}), Stern Judging (\textbf{c}) and Simple Standing (\textbf{d}), with full readout ($\alpha=\beta=1$). Dots mark pairs at which a DISC population resists both rare ALLD and rare ALLC; under every norm these are the pairs with $k_1\geq2$, together with $(1,0)$ and $(1,1)$. The most cooperative pairs are $(5,0)$ ($\bar\rho=0.888$), $(7,0)$ ($0.813$), $(7,6)$ ($0.743$) and $(2,7)$ ($0.836$), respectively. Under Simple Standing, pairs with $k_1=1$ and $k_2\geq2$ are highly cooperative but not protected, because ALLC can invade a DISC population. $b/c=5$, $e_1=e_2=0.02$.}
\label{ed:depth-evolution}
\end{figure}
\clearpage

% ED8 (working ED3; simplex panels a-d from verified/figure3_ed2-ed5_ed9/figures/ed3_simplex,
% panels e-l from verified/figure3_ed2-ed5_ed9/figures/extended_data_figure3_planes)
\begin{figure}[p]
\centering
\figureasset{figures/ED_Fig3.pdf}
\caption{\textbf{When the direction of inertia is set by readout instead of depth, tolerance again supports discriminators and cooperation, whereas strictness alone does not.}
All panels use one-step depths ($k_1=k_2=1$); columns show Scoring, Shunning, Stern Judging and Simple Standing.
\textbf{a}--\textbf{d}, Replicator dynamics, as in Fig.~\ref{fig:strategy-dynamics}, with tolerance-only readout, $(\alpha,\beta)=(1,0)$ (top), and strictness-only readout, $(\alpha,\beta)=(0,1)$ (bottom). Filled circles, equilibria stable along the boundary on which they lie; open circles, unstable equilibria. With tolerance-only readout, the DISC vertex is a strict local attractor under all four norms; with strictness-only readout it is not under any norm, and the filled ALLC--DISC equilibria under Scoring and Simple Standing can be invaded by ALLD.
\textbf{e}--\textbf{h}, Mean long-run cooperation $\bar\rho$ across the readout plane ($41\times41$ grid; 512 initial compositions per point; non-linear colour scale). Under Scoring, Shunning and Stern Judging, cooperation is highest at tolerance-only readout and vanishes at strictness-only readout; under Simple Standing it is high across the plane.
\textbf{i}--\textbf{l}, Invasion of a resident DISC population by rare ALLD and ALLC ($801\times801$ grid), with colours and curves as in Extended Data Fig.~\ref{ed:error-invasion}. A DISC population resists both mutants mainly at high tolerance strength. $b/c=5$, $e_1=e_2=0.02$.}
\label{ed:readout-evolution}
\end{figure}
\clearpage

% ED9 (working ED7; verified/figure5_ed7/figures/extended_data_figure7)
\begin{figure}[p]
\centering
\figureasset{figures/ED_Fig7.pdf}
\caption{\textbf{Protective readouts exist under all four norms, but their location depends on the norm, the allocation and the benefit-to-cost ratio.}
Resistance of a resident DISC population to the complete family of opportunistic mutants ($0<\delta\leq1$; Fig.~\ref{fig:opportunism}) for full one-step inertia and the representative multi-step allocations ($\beta=1$).
\textbf{a}, Tolerance strength $\alpha$ at $b/c=5$. Colour shows the resident reputation $r_{\mathrm R}^*$ (intended cooperation of residents); black outlines enclose the secure intervals. Open circles mark interval ends at folds of the high-reputation state, which are excluded. Multi-step intervals are $[0,0.807]$ (Scoring), $(0.591,0.806]$ (Shunning), $(0.567,0.804]$ (Stern Judging) and $[0,1]$ (Simple Standing); one-step intervals are $(0.432,0.967]$, $(0.865,0.990]$, $(0.852,0.989]$ and $[0,0.965]$.
\textbf{b}--\textbf{e}, Secure regions (black outlines) across $\alpha$ and $b/c\in[2,8]$ under Scoring (\textbf{b}), Shunning (\textbf{c}), Stern Judging (\textbf{d}) and Simple Standing (\textbf{e}) for one-step (top) and multi-step (bottom) inertia; colour, $r_{\mathrm R}^*$, which does not depend on $b/c$. The lower boundaries, set by folds, are independent of $b/c$, whereas the upper boundaries, set by $\Lambda'(0)=0$, increase with $b/c$ (Eq.~\eqref{eq:results-alpha-upper}). Circles mark $(\alpha,b/c)=(0.8,5)$, which is secure for all four multi-step allocations ($r_{\mathrm R}^*=0.978$--$0.990$) but not for one-step Shunning and Stern Judging. $e_1=e_2=0.02$.}
\label{ed:opportunism-security}
\end{figure}
\clearpage

\section*{Supplementary Information}
\textit{To be completed following the author's outline. Detailed derivations, additional validation and supplementary analyses will be added here.}
\end{document}

